% ============================================================
% Преамбула (ГОСТ 7.32-2017, XeLaTeX)
% Перенесена из курсовой работы по ТАУ; добавлен TikZ для схем
% ============================================================
\documentclass[a4paper,14pt]{extreport}

% --- Шрифты и язык (XeLaTeX) --------------------------------
\usepackage{fontspec}
\setmainfont{Times New Roman}
\usepackage{polyglossia}
\setmainlanguage{russian}
\setotherlanguage{english}
% Моноширинный шрифт с поддержкой кириллицы (для \texttt)
\setmonofont{Courier New}
\newfontfamily\cyrillicfonttt{Courier New}

% --- Геометрия страницы -------------------------------------
\usepackage[
  left   = 30mm,
  right  = 15mm,
  top    = 20mm,
  bottom = 20mm
]{geometry}

% --- Межстрочный интервал и абзацный отступ -----------------
\usepackage{setspace}
\onehalfspacing
\setlength{\parindent}{1.25cm}

% --- Математика ----------------------------------------------
\usepackage{amsmath}
\usepackage{amssymb}
% Знаки нестрогих неравенств в принятом в русской литературе начертании
\renewcommand{\le}{\leqslant}
\renewcommand{\ge}{\geqslant}
% Сквозная нумерация формул (как в образце оформления)
\counterwithout{equation}{chapter}
% ГОСТ 7.32-2017, п. 6.8.1: отбивка формулы одинакова после полной
% и неполной строки
\makeatletter
\g@addto@macro\normalsize{%
  \setlength{\abovedisplayshortskip}{\abovedisplayskip}%
  \setlength{\belowdisplayshortskip}{\belowdisplayskip}}
\makeatother
\normalsize

% --- Рисунки и таблицы --------------------------------------
\usepackage{graphicx}
\graphicspath{{figures/}}
\usepackage{float}

\usepackage{booktabs}
\usepackage{longtable}
\usepackage{array}

% --- Подписи (ГОСТ: «Рисунок N — ...», «Таблица N — ...») ---
\usepackage{caption}
\DeclareCaptionLabelFormat{gost}{#1~#2}
\DeclareCaptionLabelSeparator{gost}{ --- }
\captionsetup{
  labelsep       = gost,
  labelformat    = gost,
  justification  = centering,
  skip           = 6pt
}
\captionsetup[figure]{
  name           = Рисунок,
  position       = below,
  justification  = centering
}
\captionsetup[table]{
  name           = Таблица,
  position       = above,
  justification  = raggedright,
  singlelinecheck= false
}

% Сквозная нумерация рисунков и таблиц (как в образце оформления)
\counterwithout{figure}{chapter}
\counterwithout{table}{chapter}

% --- Стиль заголовков (titlesec, как в образце оформления) --
\usepackage{titlesec}

% Заголовок раздела: «N Название» — полужирный, кегль основного
% текста, без слова «Глава», с абзацного отступа (ГОСТ 7.32-2017, п. 6.2.1)
\titleformat{\chapter}[hang]
  {\normalfont\bfseries}{\thechapter}{1ex}{}
\titlespacing*{\chapter}{\parindent}{0pt}{14pt}

% Структурные заголовки (\chapter*: содержание, выводы,
% список литературы) — по центру, полужирные
\titleformat{name=\chapter,numberless}[block]
  {\normalfont\bfseries\filcenter}{}{0pt}{}
\titlespacing*{name=\chapter,numberless}{0pt}{0pt}{14pt}

% Подразделы — полужирные, «N.M Название» без точки после номера,
% с абзацного отступа
\titleformat{\section}[hang]
  {\normalfont\bfseries}{\thesection}{1ex}{}
\titlespacing*{\section}{\parindent}{14pt}{8pt}

\titleformat{\subsection}[hang]
  {\normalfont\bfseries}{\thesubsection}{1ex}{}
\titlespacing*{\subsection}{\parindent}{12pt}{6pt}

% Содержание (ГОСТ 7.32-2017, пп. 5.4.1, 6.13): отточие у каждой записи,
% записи разделов без полужирного и без дополнительных отбивок
\makeatletter
\renewcommand*{\l@chapter}{\@dottedtocline{0}{0pt}{1.5em}}
\makeatother

% --- Прочие пакеты ------------------------------------------
\usepackage{enumitem}
% Перечисления: с абзацного отступа, без дополнительных отбивок
% (ГОСТ 7.32-2017, п. 6.4.6)
\setlist{nosep}
\setlist[enumerate]{nosep, wide=\parindent, label=\arabic*)}
\usepackage{hyperref}
\hypersetup{
  colorlinks = true,
  linkcolor  = black,
  urlcolor   = black,
  citecolor  = black
}
% Стабилизирует файл закладок (main.out) — без него сборке не
% хватает проходов для схождения перекрёстных ссылок
\usepackage{bookmark}

% --- Схемы (TikZ) -------------------------------------------
\usepackage{tikz}
\usetikzlibrary{arrows.meta, calc, patterns}

% --- Математические сокращения ------------------------------
% Частные производные: \pd{T}{z}, \pdd{T}{z}
\newcommand{\pd}[2]{\dfrac{\partial #1}{\partial #2}}
\newcommand{\pdd}[2]{\dfrac{\partial^2 #1}{\partial #2^2}}
\DeclareMathOperator{\Div}{div}
\DeclareMathOperator{\Grad}{grad}

% Запрет «висячих» строк
\clubpenalty=10000
\widowpenalty=10000
\displaywidowpenalty=10000

% Допуск на растяжение пробелов: строки с формулами не выходят на поле
\setlength{\emergencystretch}{2em}
